\begin{lemma}[Equality of the small-radius slope limits]
Let $E$ be a finite-dimensional real inner-product space, let $D\subseteq E$ and
let $w:E\to\R$ be infinity-harmonic on $D$ in the sense of
\noderef{IsInfHarmonicOn}. Let $z\in D$. Then there is $L_*\in\R$ with
$L_*\ge0$ such that
\[
  \lim_{r\downarrow0}L_r^+(w;z)=\lim_{r\downarrow0}L_r^-(w;z)=L_* ,
\]
with the cone slopes of \noderef{ConeSlopeUpper} and
\noderef{ConeSlopeLower}.
\end{lemma}
\begin{proof}
By \noderef{IsInfHarmonicOn}, $D$ is open, so there is $r_0>0$ with
$\overline B_{2r_0}(z)\subseteq D$.

\emph{Step 1: both limits exist.}
Fix $0<r\le r'\le r_0$. Then $\overline B_{r'}(z)\subseteq D$ and
$\overline B_r(z)\subseteq D$, so \noderef{ConeComparison}, applied at the center
$z$ with radius $r$ and comparison radius $r'$, gives
\[
  L_r^+(w;z)\le L_{r'}^+(w;z),\qquad L_r^-(w;z)\le L_{r'}^-(w;z).
\]
Thus $r\mapsto L_r^\pm(w;z)$ is nondecreasing on $(0,r_0]$. Each value is
nonnegative, because the competitor $h=0$ in \noderef{ConeSlopeUpper} and
\noderef{ConeSlopeLower} contributes the quotient $0$, and each value is finite
because $w$ is continuous on the compact set $\overline B_{r_0}(z)\subseteq D$.
A bounded monotone function has a limit at the left endpoint, so
\[
  L_*^\pm:=\lim_{r\downarrow0}L_r^\pm(w;z)
       =\inf_{0<r\le r_0}L_r^\pm(w;z)
\]
exist in $\R$ and satisfy $L_*^\pm\ge0$.

\emph{Step 2: $L_*^+\le L_*^-$.}
Fix $R$ with $0<R\le r_0$. Let $0<r\le R/2$. Since $w$ is continuous on the
compact set $\overline B_r(z)$, the supremum defining $L_r^+(w;z)$ is attained:
there is $h_r$ with $|h_r|\le r$ and
\[
  L_r^+(w;z)=\frac{w(z+h_r)-w(z)}{r}.
\]
The vector $-h_r$ satisfies $|-h_r|\le r$, so it is an admissible competitor in
the supremum of \noderef{ConeSlopeLower} at the center $z+h_r$ and radius $r$,
and contributes the quotient
\[
  \frac{-\bigl(w\bigl((z+h_r)+(-h_r)\bigr)-w(z+h_r)\bigr)}{r}
   =\frac{w(z+h_r)-w(z)}{r}=L_r^+(w;z).
\]
Hence
\begin{equation}\label{eq:slopelim-cross}
  L_r^+(w;z)\le L_r^-(w;z+h_r).
\end{equation}
Also $|z+h_r-z|\le r\le R/2$, so
$\overline B_R(z+h_r)\subseteq\overline B_{2r_0}(z)\subseteq D$ and
$\overline B_r(z+h_r)\subseteq D$. Applying \noderef{ConeComparison} at the
center $z+h_r$ with radius $r$ and comparison radius $R\ge r$ gives
$L_r^-(w;z+h_r)\le L_R^-(w;z+h_r)$. Combined with
\eqref{eq:slopelim-cross},
\begin{equation}\label{eq:slopelim-cross2}
  L_r^+(w;z)\le L_R^-(w;z+h_r)\qquad(0<r\le R/2).
\end{equation}
Now let $r\downarrow0$. The left side tends to $L_*^+$ by Step 1. On the right,
$|h_r|\le r\to0$, so $z+h_r\to z$; since $\overline B_R(z)\subseteq D$, the
function $\zeta\mapsto L_R^-(w;\zeta)$ is continuous at $z$ by
\noderef{ConeSlopeContinuousInCenter}, so $L_R^-(w;z+h_r)\to L_R^-(w;z)$.
Passing to the limit in \eqref{eq:slopelim-cross2} gives
\[
  L_*^+\le L_R^-(w;z)\qquad(0<R\le r_0).
\]
Letting $R\downarrow0$ and using Step 1 gives $L_*^+\le L_*^-$.

\emph{Step 3: $L_*^-\le L_*^+$.}
The function $-w$ is infinity-harmonic on $D$. Indeed, $-w$ is continuous on the
open set $D$; and if $\phi$ is $C^2$ near $Z_0\in D$ and $(-w)-\phi$ has a local
maximum on $D$ at $Z_0$, then $w-(-\phi)$ has a local minimum on $D$ at $Z_0$, so
the supersolution inequality of \noderef{IsInfHarmonicOn} gives
$\Dinf(-\phi)(Z_0)\le0$; by \noderef{InfLap} and
$A[-\xi]=A[\xi]$ for a bilinear form $A$ we have
$\Dinf(-\phi)=-\Dinf\phi$, so $\Dinf\phi(Z_0)\ge0$, and the symmetric
computation with local minima gives the other inequality. Directly from
\noderef{ConeSlopeUpper} and
\noderef{ConeSlopeLower},
\[
  L_\rho^+(-w;z)=L_\rho^-(w;z),\qquad L_\rho^-(-w;z)=L_\rho^+(w;z)
\]
for every $\rho$, because the two defining sets of quotients are exchanged.
Applying Step 2 to $-w$ therefore gives $L_*^-\le L_*^+$.

\emph{Conclusion.} Steps 2 and 3 give $L_*^+=L_*^-$; call the common value
$L_*$, which is nonnegative by Step 1.
\end{proof}
